% Engine: XeLaTeX. No external bibliography or figures are required.
% Suggested build (not executed): xelatex -interaction=nonstopmode noted.tex
\documentclass[11pt,a4paper]{article}
\usepackage[margin=25mm,headheight=15pt]{geometry}
\usepackage{fontspec}
\IfFontExistsTF{TeX Gyre Pagella}{\setmainfont{TeX Gyre Pagella}}{\setmainfont{Latin Modern Roman}}
\IfFontExistsTF{TeX Gyre Heros}{\setsansfont{TeX Gyre Heros}}{\setsansfont{Latin Modern Sans}}
\setmonofont[Scale=MatchLowercase]{Latin Modern Mono}
\usepackage{amsmath,amssymb,mathtools,bm}
\usepackage[protrusion=true,expansion=false]{microtype}
\usepackage[table]{xcolor}
\usepackage{booktabs,tabularx,enumitem,fancyhdr}
\usepackage{hyperref}
\definecolor{ink}{HTML}{233548}
\definecolor{accent}{HTML}{216275}
\hypersetup{colorlinks=true,linkcolor=accent,urlcolor=accent,citecolor=accent,
  pdftitle={Diagnosing and stabilizing MLA latent-mean collapse},
  pdfauthor={ERDiff reproduction working notes}}
\pagestyle{fancy}
\fancyhf{}
\fancyhead[L]{\small\sffamily ERDiff reproduction}
\fancyhead[R]{\small\sffamily Source variance constraint}
\fancyfoot[C]{\small\thepage}
\setlength{\parindent}{0pt}
\setlength{\parskip}{0.45em}
\setlist[itemize]{topsep=3pt,itemsep=2pt,leftmargin=1.4em}
\setlist[enumerate]{topsep=3pt,itemsep=3pt,leftmargin=1.6em}
\DeclareMathOperator{\Var}{Var}
\newcommand{\E}{\mathbb{E}}
\newcommand{\pos}[1]{\left[#1\right]_{+}}
\newcommand{\stopgrad}{\operatorname{stopgrad}}
\allowdisplaybreaks
\begin{document}

{\sffamily\small\color{accent} RESEARCH NOTE \quad 5 NOVEMBER 2026}
\vspace{4mm}

{\LARGE\bfseries\color{ink} Diagnosing and stabilizing\\[2pt] MLA latent-mean collapse}

\vspace{2mm}
{\large A configurable source variance floor for ERDiff alignment}

\begin{quote}\small
\textbf{Status.} The stabilization option is implemented in commit
\texttt{b0a0eff}. Its tests, training experiments, and this document's XeLaTeX
build have not been run. Observations below are from user-supplied source
self-check logs; no improved result is claimed for the new constraint.
\end{quote}

\section{Question and measured evidence}

ERDiff aligns a target encoder using a frozen source diffusion model
\cite{erdiff}. Our source self-check starts from the trained source read-in,
uses the same source training trials for adaptation, and reports decoding on
the saved source validation trials. It isolates decoder degradation without
introducing a domain shift. It does not require that the paper's objective have
a stationary point at the source encoder.

All supplied runs begin at 92.16\% validation $R^2$, use model seed 42 and split
seed 42, and process 135 source training trials in shuffled batches of 45.
The following measurements are at epoch 50 of MLA-only adaptation. Variance
percentages are relative to initialization, not standard-deviation percentages.

\begin{table}[htbp]
\centering\small
\begin{tabularx}{\linewidth}{@{}Xrrr@{}}
\toprule
Training noise & Final $R^2$ & \shortstack{Mean variance\\retained} & \shortstack{Prediction variance\\retained}\\
\midrule
Uniform diffusion time & $-0.12\%$ & $0.099\%$ & $0.095\%$\\
Fixed $t=0.05$ & $74.96\%$ & $60.99\%$ & $60.75\%$\\
Fixed $t=0.2$ & $74.52\%$ & $66.03\%$ & $65.01\%$\\
Fixed $t=0.5$ & $0.41\%$ & $0.063\%$ & $0.062\%$\\
\bottomrule
\end{tabularx}
\caption{User-supplied source-only diagnostic results. The new variance
constraint was disabled in every row.}
\end{table}

The average encoder posterior variance remains approximately 99.98--99.99\%
of its initial value in these runs. Thus the dominant observed failure is
contraction of the \emph{posterior means} across observations, accompanied by
contraction of decoder outputs. It is not a global collapse of posterior
uncertainty. Aggregate statistics cannot rule out finer variance redistribution.

Fixed-noise directional finite differences mostly agree with latent autograd
derivatives within about 1\%, with some reported relative discrepancies around
2\%. These checks provide no evidence of a gross latent-gradient error in the
probed direction; they do not validate every parameter derivative. Small
directional derivatives and float32 cancellation require inspecting absolute
errors as well.

\section{What MLA optimizes}

Let $\mu_\phi(X)$ and $\ell_\phi(X)$ denote posterior mean and log variance.
Posterior sampling and VP diffusion corruption use independent Gaussian draws:
\begin{align}
 Z_0 &= \mu_\phi(X)+\exp\!\left(\tfrac12\ell_\phi(X)\right)\odot\eta,
 &\eta&\sim\mathcal N(0,I),\\
 Z_t &= a_t Z_0+b_t\epsilon,
 &\epsilon&\sim\mathcal N(0,I),\qquad b_t^2=1-a_t^2.
\end{align}
Our current MLA branch uses the trajectory-level squared residual
\begin{equation}
 \mathcal L_{\rm MLA}(\phi)
 =\E_{X,t,\eta,\epsilon}
 \left\|\epsilon-\epsilon_\theta(Z_t,t)\right\|_F^2,
 \label{eq:mla}
\end{equation}
with the source denoiser $\theta$ frozen. The paper permits time weighting;
our implementation currently uses unit weight. The beta endpoints, cutoff,
and this weighting choice are documented reproduction choices rather than
verified universal paper settings. For the current linear VP drift, its
divergence term is constant with respect to the read-in parameters and does
not change their gradient.

Minimizing a likelihood-related residual over encoder outputs does not itself
preserve distinctions between trials. An illustrative scalar Gaussian model
makes the issue explicit. Suppose the source latent has variance $s^2$ and an
ideal source denoiser predicts
\begin{equation}
 \widehat\epsilon_t(z)=k_tz,
 \qquad k_t=\frac{b_t}{a_t^2s^2+b_t^2}.
\end{equation}
If the adaptive encoding is $Z_0=cX$, where $\E X=0$ and $\Var(X)=s^2$, then
\begin{equation}
 \E\left[\epsilon-k_t(a_tcX+b_t\epsilon)\right]^2
 =(1-k_tb_t)^2+k_t^2a_t^2c^2s^2.
\end{equation}
At a nondegenerate noise level, the latter term encourages $c\to0$. Independent
fixed posterior variance adds another constant term to this toy expectation.
This illustrates mean contraction even with an ideal denoiser. It is not a
proof of the exact neural-model dynamics or a claim that the full ERDiff
objective must collapse. The paper also uses a transport regularizer.

The fixed-time results indicate that noise selection affects severity, but
both lower-noise runs still damage source decoding. Changing the noise interval
alone is therefore not an established repair.

\section{The source variance floor}

For clean source training trials, compute frozen encoder means
$M^{s}_{i\ell j}$ for trial $i$, time bin $\ell$, and latent coordinate $j$.
The source reference is the sample variance across trials at each bin:
\begin{equation}
 v^{s}_{\ell j}=\stopgrad\!\left[
 \frac{1}{N_s-1}\sum_{i=1}^{N_s}
 (M^s_{i\ell j}-\overline M^s_{\ell j})^2\right].
 \label{eq:sourcevariance}
\end{equation}
All saved source training trials contribute once, without spike augmentation
or dropped batches. The reference is computed before adapting the read-in.
It uses no source validation labels or target behavior labels.

For an adaptation minibatch of $B\ge2$ trials, let
\begin{equation}
 v^{t}_{\ell j}(\phi)=\frac{1}{B-1}\sum_{i=1}^{B}
 \left(\mu_\phi(X_i)_{\ell j}-\overline\mu_{\ell j}\right)^2.
\end{equation}
The proposed soft penalty is
\begin{equation}
 \boxed{\displaystyle
 \mathcal V(\phi)=\sum_{\ell,j}
 \pos{\rho\sqrt{v^{s}_{\ell j}+\delta}
              -\sqrt{v^{t}_{\ell j}(\phi)+\delta}}^2,}
 \qquad 0<\rho\le1,
 \label{eq:floor}
\end{equation}
where $[u]_+=\max(u,0)$ and $\delta=10^{-8}$ provides numerical smoothing.
The coefficient $\rho$ is a \emph{standard-deviation} fraction: $\rho=0.9$
corresponds approximately to an 81\% variance floor when variances are well
above $\delta$. It is a soft penalty, not a hard bound.

The default \texttt{batch} scope uses the current MLA minibatch for $B$.
The optional \texttt{full\_train} scope instead sets $B$ to the number of
\emph{unique clean adaptation training trials}. It recomputes all their means
with the current encoder at every optimizer step, retains their computation
graph, and backpropagates the exact full-pool penalty. It does not use a
detached or stale estimate, augment the pool, or include validation/test
trials. MLA itself still uses the balanced shuffled minibatch. This removes
minibatch sampling variation from the variance estimate, but the estimate
still changes as the read-in changes. The extra forward graph increases
memory and computation; this implementation does not silently approximate it
with detached chunk statistics. The penalty is applied once per optimizer
step with the same weight in either scope.

Variance is taken across trials separately at every time bin. A shared,
time-varying trajectory repeated for every trial therefore cannot satisfy
this floor merely by retaining temporal variation. Conversely, penalizing only
the per-sample posterior variance $\exp(\ell_\phi)$ would miss the dominant
failure observed in the supplied logs.

The total optional objective is
\begin{equation}
 \mathcal L=(1-\alpha)
       \left(\mathcal L_{\rm MLA}+\lambda_v\mathcal V\right)
       +\alpha\mathcal L_{\rm OT}
       +\lambda_{\rm smooth}\mathcal L_{\rm smooth}.
 \label{eq:total}
\end{equation}
Thus $\lambda_v=0$ recovers the existing baseline and $\alpha=1$ turns off the
MLA branch, including its stabilization. Both MLA and $\mathcal V$ sum over
time and latent coordinates, making their reduction conventions explicit.
Their gradient magnitudes still depend on latent scale and require measurement.

\section{Gradient and limitations}

For an active deficit
$g_{\ell j}=\rho\sqrt{v^s_{\ell j}+\delta}-\sqrt{v^t_{\ell j}+\delta}>0$,
the derivative with respect to a minibatch posterior mean is
\begin{equation}
 \frac{\partial\mathcal V}{\partial\mu_{i\ell j}}
 =-\frac{2g_{\ell j}}{(B-1)\sqrt{v^t_{\ell j}+\delta}}
       (\mu_{i\ell j}-\overline\mu_{\ell j}).
\end{equation}
Negative-gradient motion pushes distinct means away from their minibatch
mean when spread is below the reference floor. The encoder couples these
coordinates, so an actual read-in update need not follow every independent
latent direction. The frozen reference receives no gradient.

Several limits are material:
\begin{itemize}
 \item At exact collapse, all centered means are zero and this derivative can
 vanish. Start each comparison from the original source checkpoint, not an
 already collapsed read-in.
 \item Sample variance is unbiased, but its square root and the hinge are not.
 Minibatches can activate the penalty even at a well-aligned source encoder.
 Balanced batches reduce the small-tail problem, not all sampling noise.
 \item The floor constrains diagonal spread, not cross-coordinate covariance,
 trial identity, behavioral orientation, means, or decoding accuracy.
 \item A target domain may legitimately have less variation than the source.
 A strong floor can then manufacture unnecessary spread or oppose useful
 adaptation. Fixed cue-relative timing also limits the validity of a
 per-time-bin reference when movement timing differs.
 \item A finite penalty weight does not guarantee the floor is attained.
 Preserving source decoding by preventing all adaptation is not sufficient:
 recovery of controlled synthetic shifts is a separate acceptance gate.
\end{itemize}

This is an \textbf{experimental stabilization variant}, not a term asserted
to appear in the published ERDiff objective. Source self-checks identify a
failure mode; they do not establish a research contribution or paper-level
performance reproduction.

\subsection{Source-calibrated gradient scaling}

Optional gradient diagnostics measure each loss component with respect to the
target read-in parameters. For component $k$, the report records
$\|\nabla_\theta \mathcal L_k\|_2$, the norm after its current scalar
coefficient, and pairwise gradient cosines. These measurements use
\texttt{autograd.grad}; they do not populate parameter gradients or change the
optimizer state. Loss magnitudes alone do not determine which term dominates
an update.

The optional \texttt{source\_calibrated} mode estimates a single OT scale on
fixed, clean source-training probes. It makes an independent copy of the
source VAE/read-in, keeps the denoiser frozen, and uses locally seeded
posterior draws, diffusion noise, and time values. No target behavior labels
enter the calibration. For source probe $k$, define the coefficient-weighted
gradient magnitudes
\begin{equation}
 G_{\rm MLA}=\sqrt{\frac{1}{K}\sum_{k=1}^{K}
       \left((1-\alpha)\|\nabla_\theta(\mathcal L_{{\rm MLA},k}
                                  +\lambda_v\mathcal V_k)\|_2\right)^2},
 \qquad
 G_{\rm OT}=\sqrt{\frac{1}{K}\sum_{k=1}^{K}
       \left(\alpha\|\nabla_\theta\mathcal L_{{\rm OT},k}\|_2\right)^2}.
\end{equation}
When both RMS values are finite and exceed $10^{-12}$ in the calibration
dtype, the stored factors are
\begin{equation}
 c_{\rm MLA}=1,\qquad
 c_{\rm OT}=\operatorname{clip}\!\left(\frac{G_{\rm MLA}}{G_{\rm OT}},
                                      0.01,100\right).
\end{equation}
The calibrated objective is
\begin{equation}
 \mathcal L=(1-\alpha)(\mathcal L_{\rm MLA}+\lambda_v\mathcal V)
       +\alpha c_{\rm OT}\mathcal L_{\rm OT}
       +\lambda_{\rm smooth}\mathcal L_{\rm smooth}.
\end{equation}
The variance term participates only inside the MLA branch when its configured
weight is active; it is not scaled separately by $c_{\rm OT}$. An inactive
variance gradient does not produce an independent scale. Calibration is
marked invalid if either weighted component has
an effectively zero or non-finite RMS gradient; the code does not substitute a
factor in that case. A JSON artifact binds the exact factor to the source
checkpoint/data hashes, probe trial IDs, diffusion schedule, transport
settings, and fixed-fixture identity so the same factor can be reused across
target sessions. This scaling is an experimental optimization variant. It
changes relative step magnitudes, cannot correct opposed gradient directions
or semantic ambiguity, and does not establish improved decoding.

\section{Implementation and controlled comparison}

The controlled-run launcher accepts an immutable reference manifest, a protocol,
and one case. It writes the resolved case manifest before dispatching training.
All requested comparisons use model seed 42, split seed 42, and the reference's
150-epoch budget. The source-only calibration seed, mapping seed and fixed-bank
seed are independently recorded as 42; their generator namespaces remain
separate. Optimizer updates and adaptation sample exposures are recorded at
runtime because sessions can have different cohort sizes. Equal epochs alone
do not imply equal sample exposure.

\begin{table}[htbp]
\centering\small
\begin{tabularx}{\linewidth}{@{}lX@{}}
\toprule
Case & Explicit Phase 2 configuration change relative to the reference\\
\midrule
\texttt{reference} & Replay the reference settings and saved splits (R150).\\
\texttt{control} & C0: variance weight 10, standard-deviation floor 0.9, full clean training-pool scope; resampled noise, native input width, and no gradient scaling.\\
\texttt{fixed\_noise} & C0 plus a trial-indexed fixed Gaussian bank.\\
\texttt{matched\_size} & C0 plus deterministic input-width mapping to the source width.\\
\texttt{gradient\_balance} & C0 plus source-only calibrated OT gradient scaling and component-gradient logging.\\
\texttt{all} & C0 plus fixed noise, width mapping, and source-calibrated gradient scaling.\\
\bottomrule
\end{tabularx}
\caption{Predeclared controlled cases. They are comparisons, not a claim that
the combined case is better.}
\end{table}

Fixed noise changes repeated-run variance, not information leakage: the
alignment Gaussian bank is generated from independent stream namespaces and is
keyed by trial identity. The diagnostic probe noise is separate from fitting
noise. Gradient scaling is calibrated on fixed clean source-training probes;
no target behavior labels enter optimization or calibration. Every Phase 2
case keeps the denoiser frozen and verifies its state. The frozen check can run
at the end or every epoch, but it cannot be disabled.

For source-self runs, validation trials are used only for reporting, while
adaptation uses the source training split. The equal-width source self-copy is
an identity control; it does not establish recovery under unequal input widths.
The synthetic protocol applies fixed permutations, gains, masks, and actual
channel removal to disjoint source-validation adaptation and evaluation groups.
The source test split remains reserved. A channel-removal oracle is available
only as a known synthetic calibration; it is not an initialization for fitting.

Cross-day comparisons require the reference manifest to include each named
session's exact target split indices, raw-window intervals, and data SHA-256.
The launcher hashes each session file before alignment and rejects any
mismatch. It does not infer a replacement test set or silently retrain Phase
1 when a checkpoint was specified. Historical manifests with incomplete
fields may be run only as labelled known-configuration replays; they are not
promoted to exact reproductions. The output records per-session update counts,
sample exposures, channel widths, and evaluation metrics.

To create this parent manifest, first run the R150 cross-day reference into a
fresh checkpoint directory, then use
\path{build_cross_day_reference.py} to aggregate those session checkpoints.
Phase 1 alone cannot supply the held-out target partitions needed for a
controlled cross-day replay. The exporter rejects mixed settings, source
checkpoints, data fingerprints, or code revisions rather than inferring a
reference from incomplete artifacts.

The reference, variance control, noise, width, and gradient settings are
written to each resolved manifest before training begins. Case output folders
are unique and existing folders are rejected, so prior results are not
overwritten. A fixed bank and source calibration artifact are shared across
paired cases where their provenance is compatible. Native unequal widths
remain the default; matching width can discard channels and is an explicit
information-loss variant. For equal-width source self-checks it reduces to an
identity map, so the synthetic channel-removal protocol is the relevant width
recovery check.

\begin{tabularx}{\linewidth}{@{}lX@{}}
\toprule
Control & Meaning\\
\midrule
\texttt{--mla\_variance\_weight} & $\lambda_v$, nonnegative; default 0.\\
\texttt{--mla\_variance\_floor} & $\rho$, in $(0,1]$; default 0.9.\\
\texttt{--mla\_variance\_scope} & \texttt{batch} or \texttt{full\_train}; default \texttt{batch}.\\
\texttt{--alpha} & MLA/OT mixing weight; use 0 to isolate this experiment.\\
\bottomrule
\end{tabularx}

\path{mla_stabilization.py} implements Eqs.~\eqref{eq:sourcevariance} and
\eqref{eq:floor}. \path{train_phase2.py} adds the weighted penalty, reports
its raw and weighted values, and saves the reference and configuration.
Cross-day orchestration forwards corresponding optional flags. The existing
fixed-noise MLA probes continue to measure the \emph{unmodified MLA component};
they are not finite-difference checks of the new combined objective.
The diagnostic variance ratios use population variance, whereas the constraint
uses sample variance; this small denominator difference should be retained
when reading detailed numerical values.

The first planned comparison holds the following fixed: the same source
checkpoint, source training/validation indices, model seed 42, split seed 42,
probe seed 42, balanced batch policy, uniform source diffusion-time interval,
Adam learning rate 0.001, and 50 epochs. Compare
\begin{equation}
 \lambda_v\in\{0,1,10\},\qquad\rho=0.9,\qquad\alpha=0.
\end{equation}
These are predetermined source-only diagnostic candidates, not validated
hyperparameters. Each run restarts from the source checkpoint and writes to
a separate directory. Examine posterior-mean trial variance, prediction
variance, source validation $R^2$, penalty contribution, and clipping, including
early epochs. Then require a known recoverable synthetic-domain shift before
freezing settings for later target evaluation. Previously inspected target
test scores must not select these candidates.

\textbf{Verification boundary.} Focused tests were written for a zero deficit
at matching source means, an outward deficit gradient, rejection of a one-trial
batch, and detection of identical trials with a varying temporal template.
They were not executed. No new training result, numerical validation of the
constraint, typeset PDF, or performance improvement is claimed in this note.

\section{Source-only synthetic development checks}

\path{source_synthetic_checks.py} now provides a separate development harness.
It reads one source checkpoint and its saved source training/validation split.
The source training trials provide the frozen posterior bank and variance
reference. Saved validation windows are partitioned into adaptation and
evaluation groups with a local split seed, preserving connected groups of
overlapping raw intervals. Adaptation receives transformed spikes only;
evaluation uses the same transformation on the disjoint evaluation group and
reports final, fixed-duration decoding scores. Source test trials and real
target sessions are not evaluated or used for model selection.

Each transformation is fixed by \texttt{--shift\_seed} and shared across
weights 1 and 10. Identity is an additional no-shift control. With observations
written as column vectors, permutation and gain shifts take the forms
$x'=Px$ and $x'=Dx$, respectively. They admit exact first-layer inverses
$W'=WP^\top$ and $W'=WD^{-1}$, with unchanged bias and frozen remaining
layers. Gain factors are independently fixed in $[0.5,1.5)$.
The harness checks these oracle read-ins against clean source predictions;
\emph{oracle weights are never used to initialize adaptation}. Actual
adaptation starts from the ordinary source read-in for every condition.

The removal condition zero-masks a fixed approximately 20\% of channels and
retains the input width. It is lossy and has no claimed exact inverse. This
is a controlled information-loss check, not a reproduction of a session with
a physically different channel count. Every run resets the model, optimizer,
sampler and noise seed, uses a fixed final epoch, and records update counts.
The full-training penalty pool consists solely of synthetic adaptation
trials; the source reference remains the original source training bank.

The harness saves transformation tensors, development indices, raw window
intervals, checkpoint/data SHA256 fingerprints, per-run histories and read-ins,
and a summary CSV with paired decoding and variance statistics. RMSE is
reported after inversion of the saved source transform, in raw velocity
units. The clean-source score is measured on the same development evaluation
cohort; it need not equal the source checkpoint's score on all validation
trials. These validation trials already contributed to source checkpoint
selection, so this is development evidence, not untouched final confirmation.
Settings must be frozen before later source-test or real-target confirmation.
Neither the synthetic checks nor the full-pool implementation tests have
been executed during this implementation; no new result is asserted.

\section{Restoring the historical author transport objective}

The pinned 2023 author implementation, \path{model_functions/MLA_Model.py},
lines 190--195, forms squared distances between flattened sampled latent
trajectories and calls POT's \texttt{sinkhorn2}. To reproduce that behavior,
we now compute
\begin{align}
 C_{ij} &= \lVert z_i^s-z_j^t\rVert_2^2+10^{-4},
 & \bar C &= C/\max_{ij}C_{ij},\\
 \gamma_\varepsilon &\in \arg\min_{\gamma\in\Pi(a,b)}
 \langle\gamma,\bar C\rangle+
 \varepsilon\sum_{ij}\gamma_{ij}\log\gamma_{ij},
 & L_{\mathrm{OT}} &= \langle\gamma_\varepsilon,\bar C\rangle.
\end{align}
The masses $a$ and $b$ are uniform over each domain's batch; their batch
sizes may differ. The author default is $\varepsilon=0.01$. The local maximum
iteration default is now 1000, with Sinkhorn--Knopp and tolerance $10^{-9}$.
POT's installed version is recorded because the historical requirements
did not pin it.

This returned value excludes the entropy contribution, whereas the paper's
written objective includes it. We explicitly choose author-code fidelity
for this change. Autograd now traverses both the transport solver and
the maximum normalization. In particular, a detached-plan envelope gradient
for the entropy-inclusive value would not be equivalent to this gradient.
The previous unnormalized-cost setting $31.2$ is not the regularization
setting for this normalized objective. Old calibration files and controlled
references must be regenerated. No performance gain is established by this
change; the value, gradient and finite-difference tests are written but unrun.

\subsection{Numerical revision after iteration-zero transport failures}

The supplied cross-day log reports a POT numerical error at iteration zero,
followed by nonfinite alignment losses. Source-only calibration had finite
gradients ($c_{\mathrm{OT}}=56.8831$), but this does not certify target-batch
numerics. Ordinary Sinkhorn forms $K=\exp(-\bar C/\varepsilon)$; at
$\varepsilon=0.01$ its entries can approach $\exp(-100)$, making float32
reciprocal updates and their derivatives fragile. The first nonfinite tensor
was not captured in that run, so the precise initiating operation remains
unresolved.

We retain \texttt{ot.sinkhorn2}, the normalized cost, $\varepsilon=0.01$,
and the entropy-excluding return value. The numerical implementation now uses
\texttt{method="sinkhorn\_log"} with float64 inputs and returns the scalar in
the target model's dtype. Differentiable casts preserve the gradient path.
This targets the same mathematical objective, while finite-iteration values
and gradients need not be identical to a conventional solve. The numerical
method and precision are recorded in transport provenance, invalidating old
calibration artifacts. CPU and optional CUDA regression checks are written
for separated float32 domains, but remain unrun.

\subsection{Stopping at the first nonfinite operation}

Alignment now checks the source posterior bank and scales, target input,
encoder means and log variances, posterior scales and samples, MLA inputs,
noised latents and denoiser outputs, transport inputs and costs, and each
configured loss. When component-gradient logging is enabled, nonfinite
component norms also stop the run. Before any optimizer step, every available
read-in gradient must be finite and clipping uses
\texttt{error\_if\_nonfinite=True}; updated read-in parameters are checked
as well. A failure aborts the run instead of silently skipping batches or
changing the objective, and prevents the normal completed-checkpoint path.

The JSON failure record identifies the stage, epoch, batch and trial IDs,
plus tensor shapes, dtypes, devices, nonfinite counts and finite ranges.
No target behavior labels enter this check. Both evaluation entry points
reject nonfinite predictions rather than treating NaN metrics as a benchmark.
The historical failed outputs remain invalid; this revision does not
repair their read-ins. Checks for a finite forward loss with an infinite
backward derivative, optimizer non-mutation on failure, and corrupted
evaluation read-ins are written but unrun.

\section{Learned posterior variance without output clamps}

The author VAE computes log variance with a linear head without bounding
its output (\path{model_functions/VAE.py}, line 105 in the pinned
snapshot). New local Phase 1 models now use
\begin{equation}
 \ell_\phi(x)=W_\ell h_\phi(x)+b_\ell,\qquad
 q_\phi(z\mid x)=\mathcal N\!\left(\mu_\phi(x),
       \operatorname{diag}(e^{\ell_\phi(x)})\right),\qquad
 z=\mu_\phi(x)+e^{\ell_\phi(x)/2}\odot\eta,
 \quad\eta\sim\mathcal N(0,I).
\end{equation}
Previously, clamping $\ell$ to $[-5,0]$ restricted posterior variance to
$[e^{-5},1]$ and set the head's gradient to zero outside those bounds.
For the unclamped standard-normal KL,
\begin{equation}
 \frac{\partial\operatorname{KL}(q\Vert\mathcal N(0,I))}
      {\partial\ell_j}=\tfrac12(e^{\ell_j}-1),
\end{equation}
before the configured averaging and loss coefficient. This permits the
head to learn small or large uncertainty rather than saturating at the
artificial boundaries. Exponentiation can still overflow for extreme
outputs; removing clamps is not a numerical stability guarantee. We retain
the previous bias $-3$ and Xavier gain $0.1$ initialization as an explicitly
local choice, without claiming the full author architecture is reproduced.

The optional \texttt{legacy\_clamped} setting preserves previous experiments.
New checkpoints record the mode and bounds. All source-model loaders restore
that setting; checkpoints predating the field retain their original clamp.
New-format checkpoints with missing or contradictory metadata are rejected.
It would be inconsistent to remove the clamp only in Phase 2 while retaining
a denoiser trained on the clamped source posterior. Thus adopting the new
posterior requires cooperative Phase 1 retraining, after which diffusion
remains frozen throughout alignment. The source mean-trajectory variance
floor is a separate experimental constraint and remains configurable.

Tests are written for outputs and KL gradients beyond the old bounds,
posterior sampling scale, legacy and unclamped checkpoint restoration, and
incompatible metadata. They are unrun by instruction; no training or
typesetting was performed and no decoding improvement is claimed.

\section{Author-code denoising loss}

The pinned \path{model_functions/Diffusion.py}, lines 218--225, trains with
element-mean absolute noise-prediction error. Phase 1 now defaults to
\begin{equation}
 L_{\mathrm{DM}}=\frac{1}{BLd}\sum_{b,\ell,j}
       |\widehat\epsilon_{b\ell j}-\epsilon_{b\ell j}|.
\end{equation}
The previous element-mean MSE remains a named legacy option. These objectives
are not related by a constant scale: their derivatives with respect to the
predicted noise are respectively $\operatorname{sign}(\widehat\epsilon-
\epsilon)/(BLd)$ (zero at equality under PyTorch's convention) and
$2(\widehat\epsilon-\epsilon)/(BLd)$. A new Phase 1 run is required to obtain
an L1-trained denoiser. Loss identity and reduction are saved in its checkpoint.
The source latent sampling choice, continuous VP schedule and architecture
are retained and still differ from parts of the historical author pipeline.
Hand-derived value/gradient and duplicate-trial reduction checks are written
but unrun.

\begin{thebibliography}{9}\small
\bibitem{erdiff}
Y. Wang, Z. Wu, C. Li, and A. Wu.
\emph{Extraction and Recovery of Spatio-Temporal Structure in Latent Dynamics
Alignment with Diffusion Models}. NeurIPS, 2023.
\href{https://proceedings.neurips.cc/paper_files/paper/2023/file/7abbcb05a5d55157ede410bb718e32d7-Paper-Conference.pdf}{Published paper},
especially Eqs. 5, 26--28 and Algorithms 1--2.

\bibitem{authorcode}
ERDiff author implementation, pinned 2023 snapshot
\texttt{ae49f74e7fee70a2e3ea8205bf54d937ec552103}.
\url{https://github.com/yulewang97/ERDiff/tree/ae49f74e7fee70a2e3ea8205bf54d937ec552103}.
The paper and this historical code are distinct references; later releases
are not mixed into a single reproduction configuration.
\end{thebibliography}
\end{document}
